\section{The Online LCM Pipeline}
\label{sec:lcm_pipeline}
Translating raw sensor data from a running bearing into a lubrication condition estimate and, ultimately, a maintenance decision requires a sequence of distinct operations. Figure~\ref{fig:lcm_flow} shows the four-step online LCM pipeline that organises this review. Each step is described below; the scope of this review is steps 1--3, with signal processing as the primary focus.

\subsection{Pipeline Steps}
\label{subsec:flow}

An LCM system can be broken down into four steps, shown in Figure~\ref{fig:lcm_flow}:
\begin{enumerate}
    \item \textbf{Sensing}: collecting data from a bearing in operation.
    \item \textbf{Signal processing}: turning the raw sensor data into a more compact or more informative form, ready for the next step.
    \item \textbf{Lubrication condition assessment}: estimating the lubrication regime, film thickness, degradation, or contamination level from the processed data (see Section~\ref{sec:lubrication_physics}).
    \item \textbf{Maintenance decision}: deciding when and how much to relubricate, based on that assessment.
\end{enumerate}
\todo{Update in the end with the real flow and make sure that all references to this make sense.}

\begin{figure}[!h]
\centering
\begin{tikzpicture}[
  font=\sffamily,
  stage/.style={rectangle, rounded corners=5pt, draw=black!20, line width=0.5pt,
                fill=white, align=center, font=\sffamily\footnotesize,
                text width=2.2cm, minimum width=2.4cm, minimum height=1.5cm, inner sep=2pt,
                preaction={fill=black!15, transform canvas={xshift=0.8pt, yshift=-1pt}}},
  asset/.style={stage, draw=black!85, fill=black!85, text=white, text width=1.8cm,
                minimum width=2.0cm},
  arr/.style={-{Stealth[length=5pt, width=4pt]}, line width=0.9pt, black!55},
  region/.style={rounded corners=9pt},
  tag/.style={anchor=north west, font=\sffamily\scriptsize\bfseries, text=black!60,
              inner sep=0pt},
]
% Regions
\fill[region, black!4] (2.35,-1.30) rectangle (14.70,1.95);
\node[tag] at (2.55,1.77) {LCM + MAINTENANCE DECISION};
\fill[region, black!9] (2.65,-1.00) rectangle (11.65,1.45);
\node[tag] at (2.85,1.27) {LCM};

% Stages
\node[asset] (B)  at ( 1.00,0) {Ball bearing};
\node[stage] (S)  at ( 4.15,0) {{\scriptsize\bfseries\color{black!45}STEP 1}\\[1pt] Sensing};
\node[stage] (SP) at ( 7.15,0) {{\scriptsize\bfseries\color{black!45}STEP 2}\\[1pt] Signal processing and modelling};
\node[stage] (LA) at (10.15,0) {{\scriptsize\bfseries\color{black!45}STEP 3}\\[1pt] Lubrication condition assessment};
\node[stage] (MD) at (13.20,0) {{\scriptsize\bfseries\color{black!45}STEP 4}\\[1pt] Maintenance decision};

% Flow
\draw[arr] (B)  -- (S);
\draw[arr] (S)  -- (SP);
\draw[arr] (SP) -- (LA);
\draw[arr] (LA) -- (MD);
\end{tikzpicture}
\caption{Steps in online LCM of a running ball bearing. Signal processing sits between the raw sensor data and the lubrication condition assessment.}
\label{fig:lcm_flow}
\end{figure}

A given LCM system may stop at step 3, leaving the maintenance decision (step 4) as a separate downstream task.
